\documentclass[11pt,a4paper]{article}

% --- Packages ---
\usepackage[utf8]{inputenc}
\usepackage[T1]{fontenc}
\usepackage[french]{babel}
\usepackage[margin=2.2cm]{geometry}
\usepackage{xcolor}
\usepackage{listings}
\usepackage{graphicx}
\usepackage{hyperref}
\usepackage{fancyhdr}
\usepackage{tcolorbox}
\usepackage{booktabs}
\usepackage{titlesec}

% --- Couleurs Personnalisées ---
\definecolor{primaryblue}{RGB}{0, 75, 140}
\definecolor{sparkorange}{RGB}{226, 86, 0}
\definecolor{codebg}{RGB}{248, 249, 250}
\definecolor{codeframe}{RGB}{220, 224, 230}
\definecolor{commentgreen}{RGB}{46, 139, 87}
\definecolor{keywordblue}{RGB}{0, 0, 205}
\definecolor{stringred}{RGB}{178, 34, 34}

% --- Configuration Listings (Java & Bash) ---
\lstdefinestyle{customjava}{
    language=Java,
    backgroundcolor=\color{codebg},
    basicstyle=\ttfamily\footnotesize,
    keywordstyle=\color{keywordblue}\bfseries,
    stringstyle=\color{stringred},
    commentstyle=\color{commentgreen}\itshape,
    numbers=left,
    numberstyle=\tiny\color{gray},
    stepnumber=1,
    numbersep=8pt,
    tabsize=2,
    showspaces=false,
    showstringspaces=false,
    breaklines=true,
    breakatwhitespace=true,
    frame=single,
    rulecolor=\color{codeframe}
}

\lstdefinestyle{custombash}{
    language=bash,
    backgroundcolor=\color{codebg},
    basicstyle=\ttfamily\footnotesize,
    keywordstyle=\color{primaryblue}\bfseries,
    commentstyle=\color{commentgreen}\itshape,
    frame=single,
    rulecolor=\color{codeframe},
    breaklines=true
}

% --- En-tête et Pied de page ---
\pagestyle{fancy}
\fancyhf{}
\fancyhead[L]{\textcolor{primaryblue}{\textbf{TP 3 : Apache Spark \& Traitement Distribué}}}
\fancyhead[R]{\thepage}
\fancyfoot[C]{\footnotesize Big Data Frameworks -- Master / 5\textsuperscript{ème} Année}

\hypersetup{
    colorlinks=true,
    linkcolor=primaryblue,
    urlcolor=sparkorange,
    citecolor=primaryblue
}

% --- Titre du Document ---
\title{
    \vspace{-1.5cm}
    \begin{tcolorbox}[colback=primaryblue,colframe=primaryblue,arc=3mm,center]
        \centering \color{white}
        \LARGE \textbf{RAPPORT DE TRAVAUX PRATIQUES : TP 3} \\[0.3cm]
        \large \textbf{Installation d'Apache Spark, Intégration Hadoop et Applications Distribuées en Java}
    \end{tcolorbox}
    \vspace{-0.5cm}
}
\author{\textbf{Module :} Traitement Massif de Données / Big Data \\ \textbf{Architecture :} Cluster Distribué (Master M1, Workers M2 \& M3)}
\date{\today}

\begin{document}

\maketitle

\vspace{-0.5cm}
\tableofcontents
\vspace{0.5cm}
\hrule
\vspace{0.8cm}

% =========================================================================
\section{Introduction et Architecture du Cluster}
% =========================================================================

Ce TP a pour objectif le déploiement complet d'un cluster \textbf{Apache Spark 3.2.1} en mode \textit{Standalone} interconnecté avec le système de fichiers distribué \textbf{Hadoop HDFS 3.3.6}. Nous implémentons des solutions distribuées en Java exploitant à la fois l'API \textbf{RDD} (Spark Core) et l'API \textbf{DataFrame / Spark SQL}.

L'architecture matérielle virtualisée repose sur 3 nœuds en conteneurs Docker :
\begin{itemize}
    \item \textbf{M1 (hadoop-master) :} Nœud Maître hébergeant le \textit{Master Spark}, ainsi que le \textit{NameNode} et le \textit{ResourceManager} Hadoop.
    \item \textbf{M2 (hadoop-worker1) :} Nœud Esclave hébergeant un \textit{Worker Spark}, un \textit{DataNode} et un \textit{NodeManager}.
    \item \textbf{M3 (hadoop-worker2) :} Nœud Esclave hébergeant un second \textit{Worker Spark}, un \textit{DataNode} et un \textit{NodeManager}.
\end{itemize}

% =========================================================================
\section{Étapes d'Installation et Configuration}
% =========================================================================

\subsection{Étape 1 : Téléchargement et Déploiement de Spark}
Sur les trois machines (M1, M2, M3), l'archive officielle de Spark 3.2.1 a été déployée dans \texttt{/opt/spark} :
\begin{lstlisting}[style=custombash]
tar -xzf spark-3.2.1-bin-hadoop3.2.tgz -C /opt/
mv /opt/spark-3.2.1-bin-hadoop3.2 /opt/spark
\end{lstlisting}

Les variables d'environnement ont été définies dans \texttt{\textasciitilde/.bashrc} et appliquées via \texttt{source \textasciitilde/.bashrc} :
\begin{lstlisting}[style=custombash]
export SPARK_HOME=/opt/spark
export PATH=$SPARK_HOME/bin:$SPARK_HOME/sbin:$PATH
export HADOOP_CONF_DIR=$HADOOP_HOME/etc/hadoop
export SPARK_DIST_CLASSPATH=$($HADOOP_HOME/bin/hadoop classpath)
\end{lstlisting}

\subsection{Étape 2 : Configuration du Cluster (\texttt{spark-defaults.conf} et \texttt{workers})}
Sur le Master M1, nous avons configuré les paramètres de fonctionnement dans \texttt{/opt/spark/conf/spark-defaults.conf} :
\begin{lstlisting}[style=custombash]
spark.master                     spark://hadoop-master:7077
spark.submit.deployMode          client
spark.driver.memory              512m
spark.executor.memory            512m
spark.executor.cores             1
spark.hadoop.fs.defaultFS        hdfs://hadoop-master:8020
spark.driver.host                hadoop-master
spark.sql.shuffle.partitions     2
\end{lstlisting}

Le fichier \texttt{/opt/spark/conf/workers} liste les deux nœuds de travail :
\begin{lstlisting}[style=custombash]
hadoop-worker1
hadoop-worker2
\end{lstlisting}

\subsection{Étape 3 : Démarrage des Services et Validation}
Le cluster a été démarré à l'aide des scripts officiels :
\begin{itemize}
    \item Sur M1 : \texttt{\$SPARK\_HOME/sbin/start-master.sh}
    \item Sur M2 et M3 : \texttt{\$SPARK\_HOME/sbin/start-worker.sh spark://hadoop-master:7077}
\end{itemize}

La commande \texttt{jps} confirme l'activation des démons sur chaque nœud :
\begin{itemize}
    \item \textbf{M1 :} \texttt{Master}, \texttt{NameNode}, \texttt{SecondaryNameNode}, \texttt{ResourceManager}.
    \item \textbf{M2 \& M3 :} \texttt{Worker}, \texttt{DataNode}, \texttt{NodeManager}.
\end{itemize}
L'interface Web accessible sur \texttt{http://localhost:8080} affiche l'état \textbf{ALIVE} avec \textbf{2 Alive Workers}.

% =========================================================================
\section{Exercice 1 : Traitement de Texte avec Spark Core (RDD)}
% =========================================================================

\subsection{Objectifs}
L'exercice 1 requiert le développement d'une application en Java pour :
\begin{enumerate}
    \item Calculer la fréquence d'apparition de chaque mot dans un fichier texte.
    \item Dénombrer le nombre de mots de chaque longueur distincte présente dans le texte.
\end{enumerate}

\subsection{Code Source Java (\texttt{Exercice1.java})}
\begin{lstlisting}[style=customjava]
package tp3;

import org.apache.spark.SparkConf;
import org.apache.spark.api.java.JavaPairRDD;
import org.apache.spark.api.java.JavaRDD;
import org.apache.spark.api.java.JavaSparkContext;
import scala.Tuple2;

import java.util.Arrays;

public class Exercice1 {
    public static void main(String[] args) {
        if (args.length < 2) {
            System.err.println("Usage: Exercice1 <input_file> <output_dir>");
            System.exit(1);
        }

        String inputFile = args[0];
        String outputDir = args[1];

        SparkConf conf = new SparkConf().setAppName("TP3-Exercice1");
        JavaSparkContext sc = new JavaSparkContext(conf);

        // 1. Lecture du fichier texte depuis HDFS
        JavaRDD<String> lines = sc.textFile(inputFile);

        // Nettoyage et tokenisation
        JavaRDD<String> words = lines
            .flatMap(line -> Arrays.asList(line.split("\\s+")).iterator())
            .map(word -> word.replaceAll("[^a-zA-Z0-9]", "").toLowerCase().trim())
            .filter(word -> !word.isEmpty());

        // Question 1 : Frequence de chaque mot
        JavaPairRDD<String, Integer> wordCounts = words
            .mapToPair(word -> new Tuple2<>(word, 1))
            .reduceByKey(Integer::sum);

        wordCounts.saveAsTextFile(outputDir + "/word_frequency");

        // Question 2 : Nombre de mots par longueur
        JavaPairRDD<Integer, Integer> lengthCounts = words
            .mapToPair(word -> new Tuple2<>(word.length(), 1))
            .reduceByKey(Integer::sum);

        lengthCounts.saveAsTextFile(outputDir + "/word_length_count");

        sc.stop();
    }
}
\end{lstlisting}

\subsection{Exécution et Résultats HDFS}
La commande de soumission distribuée :
\begin{lstlisting}[style=custombash]
spark-submit --class tp3.Exercice1 \
  --master spark://hadoop-master:7077 \
  /opt/Lab/tp3.jar \
  hdfs://hadoop-master:8020/tp3/input/texte.txt \
  hdfs://hadoop-master:8020/tp3/output
\end{lstlisting}

\subsubsection{Résultats obtenus pour la Fréquence des Mots (\texttt{word\_frequency})}
\begin{center}
\begin{tabular}{lc|lc|lc}
\toprule
\textbf{Mot} & \textbf{Fréquence} & \textbf{Mot} & \textbf{Fréquence} & \textbf{Mot} & \textbf{Fréquence} \\
\midrule
\texttt{spark}      & 3 & \texttt{analytics}  & 1 & \texttt{scala}      & 1 \\
\texttt{apache}     & 2 & \texttt{engine}     & 1 & \texttt{python}     & 1 \\
\texttt{is}         & 2 & \texttt{unified}    & 1 & \texttt{java}       & 1 \\
\texttt{and}        & 2 & \texttt{scale}      & 1 & \texttt{provides}   & 1 \\
\texttt{data}       & 1 & \texttt{processing} & 1 & \texttt{apis}       & 1 \\
\bottomrule
\end{tabular}
\end{center}

\subsubsection{Résultats obtenus pour le Nombre de Mots par Longueur (\texttt{word\_length\_count})}
\begin{center}
\begin{tabular}{cc|cc}
\toprule
\textbf{Longueur du mot (caractères)} & \textbf{Nombre total de mots} & \textbf{Longueur du mot (caractères)} & \textbf{Nombre total de mots} \\
\midrule
1  & 2 & 6  & 4 \\
2  & 3 & 7  & 3 \\
3  & 3 & 8  & 1 \\
4  & 5 & 9  & 1 \\
5  & 7 & 10 & 1 \\
\bottomrule
\end{tabular}
\end{center}

% =========================================================================
\section{Exercice 2 : Traitement Structuré avec Spark SQL}
% =========================================================================

\subsection{Objectifs}
L'exercice 2 porte sur un ensemble de données CSV contenant la structure suivante :
\begin{center}
\texttt{(NCE, nom, prénom, spécialité, niveau, groupe, moyenne)}
\end{center}
Les questions à résoudre sont :
\begin{enumerate}
    \item Compter le nombre d'étudiants par groupe.
    \item Compter le nombre d'étudiants par spécialité.
    \item Calculer la moyenne générale par niveau et spécialité.
\end{enumerate}

\subsection{Code Source Java (\texttt{Exercice2.java})}
\begin{lstlisting}[style=customjava]
package tp3;

import org.apache.spark.sql.Dataset;
import org.apache.spark.sql.Row;
import org.apache.spark.sql.SparkSession;
import static org.apache.spark.sql.functions.avg;

public class Exercice2 {
    public static void main(String[] args) {
        if (args.length < 1) {
            System.err.println("Usage: Exercice2 <csv_file>");
            System.exit(1);
        }

        String csvPath = args[0];

        SparkSession spark = SparkSession.builder()
                .appName("TP3-Exercice2-Etudiants")
                .getOrCreate();

        // Chargement du fichier CSV avec deduction automatique des types
        Dataset<Row> df = spark.read()
                .option("header", "true")
                .option("inferSchema", "true")
                .csv(csvPath);

        System.out.println("==========================================");
        // 1. Nombre d'etudiants par groupe
        System.out.println("1. Nombre d'etudiants par groupe :");
        df.groupBy("groupe").count().show();

        // 2. Nombre d'etudiants par specialite
        System.out.println("2. Nombre d'etudiants par specialite :");
        df.groupBy("specialite").count().show();

        // 3. Moyenne generale par niveau et specialite
        System.out.println("3. Moyenne generale par niveau et specialite :");
        df.groupBy("niveau", "specialite")
          .agg(avg("moyenne").alias("moyenne_generale"))
          .show();
        System.out.println("==========================================");

        spark.stop();
    }
}
\end{lstlisting}

\subsection{Exécution et Réponses aux Questions}
La soumission du job est effectuée avec la commande :
\begin{lstlisting}[style=custombash]
spark-submit --class tp3.Exercice2 \
  --master spark://hadoop-master:7077 \
  /opt/Lab/tp3.jar \
  hdfs://hadoop-master:8020/tp3/input/etudiants.csv
\end{lstlisting}

\subsubsection{1. Nombre d'étudiants par groupe}
\begin{center}
\begin{tabular}{|c|c|}
\hline
\textbf{groupe} & \textbf{count} \\
\hline
1 & 3 \\
2 & 3 \\
\hline
\end{tabular}
\end{center}

\subsubsection{2. Nombre d'étudiants par spécialité}
\begin{center}
\begin{tabular}{|c|c|}
\hline
\textbf{specialite} & \textbf{count} \\
\hline
DSI & 3 \\
SEM & 3 \\
\hline
\end{tabular}
\end{center}

\subsubsection{3. Moyenne générale par niveau et spécialité}
\begin{center}
\begin{tabular}{|c|c|c|}
\hline
\textbf{niveau} & \textbf{specialite} & \textbf{moyenne\_generale} \\
\hline
3 & DSI & 14.0 \\
2 & SEM & 10.5 \\
3 & SEM & 13.5 \\
\hline
\end{tabular}
\end{center}

% =========================================================================
\section{Conclusion}
% =========================================================================

Ce TP a mis en évidence l'efficacité et la flexibilité d'Apache Spark dans un environnement distribué :
\begin{itemize}
    \item L'intégration native avec \textbf{Hadoop HDFS} permet un découpage automatique des blocs et une localité des données optimale.
    \item L'API bas niveau \textbf{RDD} offre un contrôle précis sur les transformations personnalisées (Word Count et longueurs de mots).
    \item L'API haut niveau \textbf{Spark SQL (DataFrame)} optimise automatiquement les plans d'exécution via le moteur Catalyst pour les agrégations relationnelles (groupements et moyennes sur les étudiants).
\end{itemize}

\end{document}

